\section{Signed Circulants and Flux Coordinates}\label{sec:preliminaries}

\subsection{Signed adjacency matrices and switching}

For an integer \(n\geq 8\), let \(G_n=C_n(1,2)\) have vertex set \(\mathbb{Z}/n\mathbb{Z}\) and edges

\begin{equation}
\{i,i+1\}, \{i,i+2\}    (i \in \mathbb{Z}/n\mathbb{Z}).
\end{equation}

A signing is a map \(\sigma:E(G_n)\to \{+1,-1\}\). Its signed adjacency matrix
\(A_\sigma\) is the real symmetric matrix whose \(ij\) entry is the sign of
\(\{i,j\}\) when this is an edge and zero otherwise. We write

\begin{equation}
\rho(A_\sigma)=\max\{\lvert \lambda\rvert: \lambda \in \operatorname{spec}(A_\sigma)\}.
\end{equation}

For a vertex-sign function \(d:\mathbb{Z}/n\mathbb{Z}\to \{+1,-1\}\), let \(D=\operatorname{diag}(d_i)\). Replacing
\(\sigma_ij\) by \(d_i \sigma_ij d_j\) replaces \(A_\sigma\) by \(D A_\sigma D\) and
therefore preserves the spectrum. This operation is switching.

Because \(G_n\) is connected and has \(2n\) edges, its cycle space has dimension
\(2n-n+1=n+1\). Hence there are exactly \(2^{n+1}\) switching classes. We use a
cycle-flux coordinate system adapted to the Hamilton cycle of step-one edges.

\subsection{Hamilton gauge, triangle flux, and quadrilateral flux}

Let \(a_i\) be the sign of the step-one edge \(\{i,i+1\}\) and \(b_i\) the sign of
the step-two edge \(\{i,i+2\}\). Define the triangle flux and Hamilton holonomy by

\begin{equation}
\begin{aligned}
\tau_i&=a_i a_{i+1} b_i, \\
\alpha&=\prod_{i=0}^{n-1} a_i.
\end{aligned}
\end{equation}

The \(n\) triangle cycles together with the step-one Hamilton cycle form a cycle
basis, so \((\tau_{0},\ldots,\tau_{n-1},\alpha)\) determines the switching class.

Switch so that all step-one edges are positive except possibly the edge
crossing the cut from \(n-1\) to \(0\). Equivalently, extend a vector to the
integer lattice with twisted boundary condition

\begin{equation}
x_{i+n}=\alpha x_i.
\end{equation}

The signed operator then has the local form

\begin{equation}
\label{eq:periodic-operator}
(A_\tau x)_i=x_{i-1}+x_{i+1}+\tau_{i-2}x_{i-2}+\tau_i x_{i+2}.
\end{equation}

The distinguished local quadrilateral flux word is

\begin{equation}
\label{eq:quadrilateral-flux}
Q_i=\tau_i \tau_{i+1}.
\end{equation}

\begin{figure}[tbp]
\centering
\begin{tikzpicture}[
  x=1.45cm,y=1.05cm,
  vertex/.style={circle,draw,fill=white,inner sep=1.8pt},
  edgeone/.style={line width=0.75pt},
  edgetwo/.style={line width=0.75pt,dashed},
  flux/.style={font=\small}
]
  \node[vertex,label=below:{\(i\)}] (v0) at (0,0) {};
  \node[vertex,label=below:{\(i+1\)}] (v1) at (1,0) {};
  \node[vertex,label=below:{\(i+2\)}] (v2) at (2,0) {};
  \node[vertex,label=below:{\(i+3\)}] (v3) at (3,0) {};

  \draw[edgeone] (v0)--node[below=3pt,flux]{\(a_i\)} (v1);
  \draw[edgeone] (v1)--node[below=3pt,flux]{\(a_{i+1}\)} (v2);
  \draw[edgeone] (v2)--node[below=3pt,flux]{\(a_{i+2}\)} (v3);
  \draw[edgetwo,bend left=38] (v0) to node[above=3pt,flux]{\(b_i\)} (v2);
  \draw[edgetwo,bend left=38] (v1) to node[above=3pt,flux]{\(b_{i+1}\)} (v3);

  \node[font=\scriptsize,align=center] at (0.72,1.05)
    {\(\tau_i=a_i a_{i+1}b_i\)};
  \node[font=\scriptsize,align=center] at (2.28,1.05)
    {\(\tau_{i+1}=a_{i+1}a_{i+2}b_{i+1}\)};
  \node[flux,draw,rounded corners=1pt,inner sep=4pt] at (1.5,1.82)
    {\(Q_i=\tau_i\tau_{i+1}\)};

  \draw[-{Latex[length=2mm]},thin] (1.20,1.61)--(0.90,1.25);
  \draw[-{Latex[length=2mm]},thin] (1.80,1.61)--(2.10,1.25);
\end{tikzpicture}

\caption{Local flux coordinates in \(C_n(1,2)\). Solid edges have cyclic
distance one and dashed arcs have distance two. The adjacent triangle fluxes
share the edge sign \(a_{i+1}\); their product is the quadrilateral flux
\(Q_i\). The drawing fixes notation only and does not impose a gauge.}
\label{fig:flux-coordinates}
\end{figure}

For a \(p\)-periodic word \(\tau\), Equation~\eqref{eq:quadrilateral-flux} implies
\(\prod_{i=0}^{p-1} Q_i=1\). Conversely, any word
\(Q \in \{+1,-1\}^p\) with product \(+1\) has exactly two periodic lifts, obtained by
choosing \(\tau_{0}=+1\) or \(-1\) and recursing via
\(\tau_{i+1}=Q_i \tau_i\).

We call

\begin{equation}
D(Q)=\{i:Q_i=+1\}
\end{equation}

the defect set and write \(d=\lvert D(Q)\rvert\). We also use the cyclic local statistics

\begin{equation}
\label{eq:local-defect-statistics}
\begin{aligned}
a&=\#\{i:Q_i=Q_{i+1}=+1\}, \\
b&=\#\{i:Q_i=Q_{i+2}=+1\}.
\end{aligned}
\end{equation}

When \(n=8\), the graph contains two additional step-two quadrilaterals not
among the local cycles encoded by \eqref{eq:quadrilateral-flux}. This causes no loss of finite
information: the cycle-basis coordinates \((\tau,\alpha)\), rather than \(Q\) alone,
parametrize every switching class and are the coordinates used in the
all-signings enumeration.

\subsection{Floquet matrices}\label{subsec:floquet-matrices}

Suppose \(\tau\) has displayed period \(p\). Write an integer index as \(mp+r\),
where \(0\leq r<p\), and seek Bloch vectors of the form

\begin{equation}
x_{mp+r}=z^m v_r.
\end{equation}

Substitution in \eqref{eq:periodic-operator} gives a \(p\times p\) Laurent matrix \(H_\tau(z)\). More
explicitly, for each transition from output residue \(r\) to an absolute source
index \(j\), write \(j=mp+s\), \(0\leq s<p\), and add the transition coefficient times
\(z^m\) to entry \((r,s)\). This convention automatically adds coefficients when
different transitions collide in short cells.

The identity

\begin{equation}
\label{eq:fiber-transpose}
H_\tau(z)^T=H_\tau(z^{-1})
\end{equation}

follows by reversing each undirected transition. Thus \(H_\tau(z)\) is Hermitian
on the unit circle.

For an infinite periodic phase, define

\begin{equation}
\label{eq:squared-periodic-radius}
R(Q)=\sup_{\lvert z\rvert=1} \rho(H_\tau(z))^{2}.
\end{equation}

The two lifts of \(Q\) have the same value, so this notation is unambiguous.
Notice that \(R(Q)\) is always a \textbf{squared} spectral radius.

For every explicit fiber certificate below we use the canonical lift
\(\tau_{0}=+1\), \(\tau_{i+1}=Q_i \tau_i\), and the ordered fiber basis
\((v_{0},\ldots,v_{p-1})\). The other lift is obtained by the diagonal conjugacy and
fiber reparametrization in Lemma~\ref{lem:operator-equivalences}, so it has the same squared spectral
conclusions.

Whenever an explicit matrix representative is required, we write
\(H_Q(z):=H_{\tau^{\mathrm{can}}}(z)\) for the fiber formed from this canonical lift and
basis. The notation \(R(Q)\) and the moments \(M_k(Q)\) remain lift-invariant.

If a finite graph has order \(n=pL\) and Hamilton holonomy \(\alpha\), its cell
sequence satisfies \(u_{m+L}=\alpha u_m\). The unitary cell shift has eigenvalues
exactly

\begin{equation}
\label{eq:finite-bloch-grid}
z^L=\alpha,
\end{equation}

and the finite matrix decomposes as

\begin{equation}
\label{eq:finite-floquet-decomposition}
A_{pL,\alpha}  \cong  \bigoplus_{z^L=\alpha} H_\tau(z).
\end{equation}

Equation~\eqref{eq:finite-bloch-grid} is a discrete finite set. Equation \(\lvert z\rvert=1\) in \eqref{eq:squared-periodic-radius} describes
the infinite-volume spectrum. We will not substitute one for the other.

\subsection{Operator equivalences}

We record the equivalences used in classifying periodic phases.

\begin{lemma}[translation, reflection, and edge-sign negation]
\label{lem:operator-equivalences}
Translating \(\tau\),
reflecting it by \(\tau_i \to  \tau_{-i-2}\), or replacing \(\tau\) by \(-\tau\) preserves
\(R(Q)\). On flux words reflection is \(Q_i \to  Q_{-i-3}\).

\end{lemma}
\begin{proof}
Let \((T_r x)_i=x_{i-r}\) and \((Jx)_i=x_{-i}\). Direct substitution in
\eqref{eq:periodic-operator} gives

\begin{equation}
\begin{aligned}
T_r A_\tau T_r^{-1}&=A_{\text{translated }\tau}, \\
J A_\tau J&=A_{\text{reflected }\tau}.
\end{aligned}
\end{equation}

Both are unitary conjugacies. For global edge-sign negation let
\((Dx)_i=(-1)^i x_i\). The
endpoints of a step-one edge have opposite \(D\) signs while those of a step-two
edge have equal signs. Therefore

\begin{equation}
A_{-\tau}=-D A_\tau D.
\end{equation}

The spectrum is negated and its square is unchanged. In an odd displayed cell
the fiber coordinate changes from \(z\) to \(-z\), but the full unit circle is
preserved.

\end{proof}
\begin{lemma}[zone folding]
\label{lem:zone-folding}
If \(\tau\) has primitive period \(q\) and is
displayed in a repeated cell \(p=mq\), then

\begin{equation}
\label{eq:zone-folding}
H_p(z) \cong \bigoplus_{w^m=z} H_q(w).
\end{equation}

In particular repetition of the unit cell preserves \(R(Q)\).

\end{lemma}
\begin{proof}
Write a residue in the repeated cell as \(r+kq\), where
\(0\leq r<q\) and \(0\leq k<m\). Decompose the \(mq\)-dimensional fiber according to the
unitary internal translation by \(q\). On its \(w\)-eigenspace vectors have the
form

\begin{equation}
x_{r+kq}=w^k v_r.
\end{equation}

The repeated boundary condition is equivalent to \(w^m=z\). Since all
coefficients in \eqref{eq:periodic-operator} are \(q\)-periodic, every transition factors out \(w^k\)
and the remaining action on \(v\) is exactly \(H_q(w)\). The \(m\) roots of
\(w^m=z\) give mutually orthogonal eigenspaces and their dimensions sum to
\(mq\). This proves \eqref{eq:zone-folding}, including multiplicities.

\end{proof}
\subsection{Closed-walk moments}

For a periodic phase define

\begin{equation}
\label{eq:moment-and-excess}
\begin{aligned}
M_k(Q)&=\operatorname{CT}_{z} \operatorname{tr}(H_Q(z)^{2k}), \\
F_k(Q)&=M_{k+1}(Q)-8M_k(Q).
\end{aligned}
\end{equation}

Constant-term extraction equals normalized integration around the unit circle:

\begin{equation}
\label{eq:moment-integral}
M_k(Q)=\frac{1}{2\pi}\int_{0}^{2\pi}\sum_j\lambda_j(e^{i\theta})^{2k}\,d\theta.
\end{equation}

It is also the signed sum of closed walks of length \(2k\) per period cell, in
the standard trace--closed-walk correspondence for graph spectra~\cite{CvetkovicRowlinsonSimic2010}.

\begin{lemma}[moment barrier]
\label{lem:moment-barrier}
If \(R(Q)\leq 8\), then
\(M_{k+1}(Q)\leq 8M_k(Q)\) for every \(k\geq 1\). Equivalently,

\begin{equation}
\label{eq:moment-barrier}
F_k(Q)>0  \Longrightarrow  R(Q)>8.
\end{equation}

\end{lemma}
\begin{proof}
Under \(R(Q)\leq 8\), every
\(y_j(\theta)=\lambda_j(e^{i \theta})^{2}\) lies in \([0,8]\). Hence
\(y_j^{k+1}\leq 8y_j^k\). Sum over \(j\) and integrate using \eqref{eq:moment-integral}.

\end{proof}
We emphasize that \eqref{eq:moment-barrier} is one-way. A nonpositive excess, or any finite list
of nonpositive excesses, does not prove \(R(Q)\leq 8\).

\subsection{Distinguished constants}

For even \(n\geq 8\), the conjectured threshold inherited from \cite{Suvagiya2026Signed} is

\begin{equation}
\label{eq:conjectured-threshold}
\begin{aligned}
\rho_-(n)&=2 \sqrt{\cos^{2}(\frac{\pi}{n})+\cos^{2}(\frac{2\pi}{n})}, \\
\rho_-(n)^{2}&=4+2\cos(\frac{2\pi}{n})+2\cos(\frac{4\pi}{n}).
\end{aligned}
\end{equation}

The target period-eight squared spectral edge is

\begin{equation}
\label{eq:target-edge}
\begin{aligned}
\eta&=4+\sqrt{10+2\sqrt{5}}, \\
\rho_*&=\sqrt{\eta}.
\end{aligned}
\end{equation}

These quantities play different roles: \(\rho_-(n)\) is the finite conjectured
lower bound, while \(\eta\) is the exact infinite-volume squared radius of the
counterexample phase.
